% ============================================================
%  Mosaic — Pasteur Labs template version.
%  Same content as ../main.tex, restyled via pasteurlabs.sty.
% ============================================================

\documentclass[11pt,a4paper]{article}
\usepackage{pasteurlabs}

\bibliographystyle{abbrvnat}

% Extras used by the body that pasteurlabs.sty does not load.
\usepackage{amsfonts}
\usepackage{nicefrac}
\usepackage{multirow}
\usepackage[capitalize,noabbrev]{cleveref}
\usepackage{orcidlink}

% Float placement: discourage float-only pages, allow more text+float
% co-placement so appendix figures stop creating mostly-blank pages.
\renewcommand{\topfraction}{0.9}
\renewcommand{\bottomfraction}{0.8}
\renewcommand{\textfraction}{0.07}
\renewcommand{\floatpagefraction}{0.85}
\setcounter{topnumber}{3}
\setcounter{bottomnumber}{2}
\setcounter{totalnumber}{5}

% ── TITLE-PAGE METADATA ─────────────────────────────────────
\title{Mosaic: A Benchmark Suite for Differentiable Physics Solvers}

\author{Andrin Rehmann\,\orcidlink{0009-0001-6002-5096},\;
        Heiko Zimmermann\,\orcidlink{0000-0001-9190-4328},\;
        Dion Häfner\,\orcidlink{0000-0002-4465-7317}}

\date{05-2026}

\organization{\includegraphics[height=26pt]{figures/pl-labs}}

\affiliation{Pasteur Labs, Brooklyn, NY, USA}

\correspondence{dion.haefner@simulation.science}

\shorttitle{Mosaic: A Benchmark Suite for Differentiable Physics Solvers --- Rehmann et al., 2026}


% ============================================================
\begin{document}
% ============================================================

\maketitle

% ===========================================================================
\begin{plabstract}
Differentiable partial differential equation (PDE) solvers underpin
solver-in-the-loop ML training, gradient-based optimal control, and
inverse problems, yet the practical cost of obtaining correct, usable
gradients from a given solver on a given problem is largely
undocumented. Integration effort, computational cost, gradient accuracy,
and numerical conditioning vary widely across solvers and are
discoverable only by trial and error.
%
We introduce \textit{Mosaic}, an extensible benchmarking framework for
differentiable PDE solvers that standardizes access to solver gradients.
Each solver is packaged as a containerized component
(\textit{Tesseract}) exposing a uniform gradient API regardless of
language or automatic differentiation (AD) strategy, enabling
researchers to evaluate, compare, and build on non-trivial physical solvers.
%
Our evaluation of 14 solvers across fluid dynamics, structural
mechanics, and heat transfer demonstrates that the benchmark surfaces
practically relevant differences: order-of-magnitude variation in
computational cost and Jacobian conditioning, alongside structural
incompatibilities that eliminate solvers from realistic tasks entirely.
Despite this variation, all solvers that
produce gradients converge to similar optima, indicating that the
practical barriers are memory limits, numerical stability, and setup
compatibility rather than gradient accuracy alone.
%
Mosaic is open-source and available at \url{https://github.com/pasteurlabs/mosaic}.
\end{plabstract}

% ---------------------------------------------------------------------------
% Teaser / visual abstract
% ---------------------------------------------------------------------------
\begin{figure}[h!]
  \centering
  \includegraphics[width=\linewidth]{figures/visual_abstract.pdf}
  \caption{Overview of Mosaic. Diverse solvers are
    wrapped behind a uniform containerized interface, enabling cross-solver comparison on shared benchmark tasks involving different physical domains.}
  \label{fig:teaser}
\end{figure}

% ===========================================================================
\section{Introduction}
\label{sec:intro}


% The problem: opaque effort, hidden incompatibilities
A growing number of physics solvers now advertise gradient support,
spanning multiple languages, AD frameworks, and numerical methods.
Yet the practical cost of obtaining correct, usable gradients from any
of these solvers is opaque. Solvers range from turnkey applications to
composable frameworks and low-level kernel toolkits, each demanding
different expertise and integration effort. Runtime overhead, gradient accuracy, and numerical conditioning
have not been systematically measured, and solver-specific constraints
(periodic-only boundary conditions,
incompatible penalization schemes, stability limits) surface only after
a practitioner has committed to a particular solver.

% Why it matters: differentiable solvers enable a broad family of techniques

Differentiable solvers enable topology
optimization~\cite{bendsoe2003topology}, aerodynamic shape
optimization~\cite{jameson1988aerodynamic}, optimal
control~\cite{bewley2001dns}, and solver-in-the-loop ML
training~\cite{kochkov2021machine, um2020sol} by backpropagating
through a numerical simulation's forward pass.
Differentiating through ODE integrators~\cite{chen2018neural,rackauckas2021universal}
and renderers~\cite{kato2020differentiable} is now routine, but PDE
solvers are more complex, less standardized, and their gradient
properties far less understood.

Learned surrogates and physics foundation models offer an alternative
path but require per-domain training data, introduce approximation
error, and currently cover only a few well-studied PDE
families~\cite{pathak2022fourcastnet, bodnar2024aurora}.
Decades of engineering have produced solvers that practitioners
trust for forward simulation. What is missing is not the solvers
themselves, which exist in abundance, but a standardized way to
access, compare, and trust their gradients.

We introduce \textit{Mosaic}, an extensible benchmark suite that
tests solver gradients in isolation, on standalone optimization
problems without learned components, where finite differences provide
ground truth and the solver's gradient mechanism is the only variable
under test. The primary contribution is the benchmark
infrastructure itself, not the empirical results from any single
instantiation. Researchers can use Mosaic to evaluate which solver fits
their problem, build solver-in-the-loop pipelines on top of the
uniform interface, and stress-test failure modes before
committing to a particular stack. Our contributions are:
%
\begin{enumerate}\itemsep2pt
  \item A \textbf{benchmark suite and containerized solver interface}
    spanning 14 solvers across four physical domains (heat transfer,
    structural mechanics, 2D and 3D Navier--Stokes;
    \cref{tab:domains}). Each solver is wrapped behind a uniform
    forward/gradient API regardless of language or AD strategy, so
    researchers can run and compare real solvers without
    solver-specific expertise.
  \item A \textbf{standardized evaluation protocol} measuring gradient
    accuracy, computational cost, Jacobian conditioning, setup
    compatibility, and end-to-end optimization convergence.
  \item \textbf{Empirical findings} that validate the benchmark by
    surfacing practically relevant differences: solver choice can
    determine whether gradient-based workflows succeed or fail,
    independent of the optimizer,
    with a negative control that isolates AD bugs from chaotic physics.
  \item An \textbf{open-source release} of all solver interfaces, evaluation
    code, and templates, extensible to new solvers and domains, with
    continuous re-evaluation designed to keep published results current
    as solvers improve.
\end{enumerate}

% ===========================================================================
\section{Background and related work}
\label{sec:related}

\paragraph{Adjoint methods and differentiable physics}
Adjoint methods for PDE-constrained optimization date back to
Pironneau~\cite{pironneau1974optimum} and
Jameson~\cite{jameson1988aerodynamic}, with discrete adjoint
formulations following in~\cite{giles2000introduction,nadarajah2000studies}.
Modern AD frameworks now differentiate solver code automatically:
JAX source transformation~\cite{bradbury2018jax}, PyTorch operator
overloading~\cite{paszke2019pytorch}, Julia source-to-source AD via
Zygote~\cite{innes2018don}, LLVM-IR-level AD via
Enzyme~\cite{moses2020enzyme}, NVIDIA
Warp kernel-level AD~\cite{macklin2024warp}, and tape-based discrete
adjoints in dolfin-adjoint and
firedrake-adjoint~\cite{farrell2013automated}.
Hand-derived adjoints are fast but tied to a fixed set of objectives,
requiring new derivations whenever the control variable or loss function
changes. Algorithmic AD removes this constraint at the price of higher
memory or computational cost, a tradeoff whose magnitude has remained
largely anecdotal. Mosaic measures it empirically across solver families.

\paragraph{Learned surrogates and physics foundation models}
Neural operator surrogates~\cite{kovachki2023neural,li2021fourier,lu2021learning},
PINNs~\cite{raissi2019physics}, and physics foundation
models~\cite{pathak2022fourcastnet,bodnar2024aurora,price2024gencast,hao2024dpot,herde2024poseidon}
are complementary to differentiable solvers but trade accuracy for speed,
degrade outside the training distribution, and require solver-generated
data in the first place. Mosaic benchmarks the solver layer directly.

\paragraph{Existing benchmarks and evaluation efforts}
Several benchmarks evaluate PDE solvers or learned surrogates, but none
systematically measures gradient properties across solver backends.
APEBench~\cite{koehler2024apebench} benchmarks neural PDE emulators on
forward accuracy without evaluating solver gradients.
PDEBench~\cite{takamoto2022pdebench} provides datasets of PDE solutions
for surrogate training but includes no gradient evaluation.
NeuralFluid~\cite{fei2024neuralfluid} benchmarks differentiable
Navier-Stokes solvers on control tasks and reports gradient runtime, yet
does not systematically assess gradient accuracy, conditioning, or setup compatibility across backends.
List et al.~\cite{list2024differentiability} analyze gradient behavior in
unrolled training, focusing on convergence dynamics rather than
cross-solver comparison.
The PhiFlow benchmark~\cite{holl2024phiflow} compares AD backends within
a single framework, not across independent solver implementations.
preCICE~\cite{bungartz2016precice} standardizes solver coupling but
provides no gradient interface.
The closest analogue in design philosophy is OpenAI
Gym~\cite{brockman2016openai}, which standardized access to RL
environments behind a uniform \texttt{step}/\texttt{reset} interface.
Gym ships no RL algorithms, yet the standardized interface enabled a
wealth of downstream research by lowering the barrier to entry.
Mosaic follows the same model for differentiable solvers: solver
backends are simulators behind a standardized gradient interface, and
the benchmark ships no ML components itself but provides the
infrastructure on which solver-in-the-loop ML research can build.

% ===========================================================================
\section{The Mosaic benchmark suite}
\label{sec:mosaic}

Mosaic organizes evaluation around three concepts. A \emph{physical
setup} is a concrete problem instance: a specific geometry, boundary
conditions, and operating regime within a physical domain (e.g.,
incompressible fluids). A \emph{benchmark task} is a fixed optimization
problem on a specific setup: a control variable, an objective, and a
reference solution. A \emph{solver backend} is
the software that evaluates the forward map and (optionally) its
gradient, independent of abstraction level. Every solver that supports a
task's physical domain can be compared on it.


\subsection{Physical domains and benchmark tasks}
\label{sec:domains}

Benchmark tasks are representative optimization problems (inflow
optimization, topology optimization, inverse problems) that require
gradients for efficient solution, are supported by multiple backends
for cross-solver comparison, and contain no learned components so solver
gradients can be validated in isolation. Tasks are intentionally simpler
than state-of-the-art challenges: complexity in the physics would
confound the gradient properties we aim to isolate, so any failure can
be attributed to the solver's gradient mechanism rather than to an
ill-posed optimization problem. The heat domain (H) additionally serves
as a negative control: its forward problem is smooth and elliptic, so
any gradient failure on H is attributable to an AD bug rather than to
chaotic dynamics. The inversion itself is poorly conditioned (its
Hessian spans about eight orders of magnitude), which affects how fast
optimizers converge rather than whether gradients are correct.

The four domains form a difficulty progression
(\cref{tab:domains,fig:domains}; full specifications in the
supplementary material).
\begin{table}[t]
  \centering\small
  \caption{Physical domains and optimization tasks. The backend count
    includes forward-only references; per-solver details are in
    \cref{tab:solvers}.}
  \label{tab:domains}
  \begin{tabularx}{\linewidth}{c>{\hsize=0.95\hsize}X>{\hsize=1.05\hsize}Xlll}
    \toprule
    \textbf{ID} & \textbf{Domain}
      & \multicolumn{2}{c}{\textbf{Optimization task}} & \textbf{Backends} \\
    \cmidrule(lr){3-4}
    & & \textit{Task} & \textit{Control dim.} & \\
    \midrule
    H & Heat transfer (slab) & Conductivity inversion
      & 128 & 5 \\
    & \multicolumn{2}{l}{\textit{%
      Negative control: smooth physics isolates AD bugs}} & & \\[5pt]
    S & Structural mech.\ (cantilever) & Compliance minimization (SIMP)
      & 2\,048 & 6 \\
    & \multicolumn{2}{l}{\textit{%
      Well-studied topology optimization with known optimal}} & & \\[5pt]
    F2 & Incompr.\ fluids (2D cylinder) & Inflow optimization for drag min.
      & 32 & 7 \\
    & \multicolumn{2}{l}{\textit{%
      Anchor domain; steady flow at $\mathrm{Re}{=}20$}} & & \\[5pt]
    F3 & 3D Navier--Stokes (TGV) & Initial condition recovery
      & 12\,288 & 7 \\
    & \multicolumn{2}{l}{\textit{%
      Scale test; periodic box isolates time-stepping gradients}} & & \\
    \bottomrule
  \end{tabularx}
\end{table}

\begin{figure}[t]
  \centering
  \begin{subfigure}[t]{0.48\linewidth}
    \centering
    \includegraphics[width=0.95\linewidth]{figures/domain4_heat.pdf}
    \caption{\textbf{H}\,: Steady-state heat conduction --- conductivity inversion.}
    \label{fig:domain_h}
  \end{subfigure}\hfill
  \begin{subfigure}[t]{0.48\linewidth}
    \centering
    \includegraphics[width=0.95\linewidth]{figures/domain3_structures.pdf}
    \caption{\textbf{S}\,: Cantilever beam --- compliance minimization (SIMP).}
    \label{fig:domain_s}
  \end{subfigure}\\[0.2em]
  \begin{subfigure}[t]{0.48\linewidth}
    \centering
    \includegraphics[width=0.95\linewidth]{figures/domain1_2d_fluids.pdf}
    \caption{\textbf{F2}\,: 2D cylinder flow --- inflow optimization.}
    \label{fig:domain_f2}
  \end{subfigure}\hfill
  \begin{subfigure}[t]{0.48\linewidth}
    \centering
    \includegraphics[width=.85\linewidth]{figures/domain2a_3d_ic_recovery.pdf}
    \caption{\textbf{F3}\,: 3D periodic box --- initial condition recovery.}
    \label{fig:domain_f3}
  \end{subfigure}
  \caption{Benchmark tasks across the four physical domains in Mosaic,
    ordered by difficulty.
    Each panel shows the control variable, the physical process,
    and the optimization objective.}
  \label{fig:domains}
\end{figure}

\subsection{Design principles}
\label{sec:scope}

Mosaic consists of three artifacts: a containerized solver interface,
an evaluation protocol, and a library of benchmark task templates.
Containerization is handled by Tesseract~\cite{haefner2025tesseract},
an existing open-source runtime for packaging scientific codes; Mosaic
contributes the evaluation methodology, benchmark tasks, and empirical
results on top of it. The four domains and 14 solver backends described
in this paper are the initial instantiation; the reusable evaluation
infrastructure is the lasting contribution, designed to evolve as the
community adds solvers, improves configurations, and extends to new
domains.

\paragraph{Solver interface and extensibility}
Each solver is a containerized component exposing two entry points,
\texttt{apply(inputs)\,$\rightarrow$\,outputs} and
\texttt{vector\_jacobian\_product(inputs, cotangents)\,$\rightarrow$\,vjp}.
The container isolates all dependencies (language, libraries, GPU
runtime), so any of the 14 backends can be swapped by changing a single
configuration string.
Forward-only solvers
(OpenFOAM, deal.II) implement \texttt{apply} only; Mosaic provides
finite-difference gradients automatically at $O(N)$ cost, serving as the
reference baseline.
Because solver APIs differ too widely for a single generic wrapper, each
benchmark task defines a \emph{task-specific interface schema} (inputs,
outputs, shapes, and units) that multiple solvers implement. Adding a new
backend requires only implementing the task schema and adding a container;
the evaluation suite re-runs automatically on every proposed change.

Benchmark tasks are templates, not fixed endpoints: each template is a
validated starting configuration with known-good settings, so users can
modify the physical setup (geometry, Reynolds number, boundary conditions)
while inheriting the evaluation protocol and finite-difference reference.
Details on contributing and containerization overhead are in \cref{app:software}.

Integration effort per backend is quantified in \cref{tab:merged}; column
definitions and full experimental specifications are in \cref{app:effort}.

\subsection{Evaluation protocol}
\label{sec:protocol}

All metrics are computed uniformly across solvers. The protocol is
solver-agnostic: it operates on the containerized interface (forward and
VJP) and never inspects solver internals.

\begin{description}\itemsep2pt
  \item[Setup compatibility.] For each solver--task pair, we record
    whether the solver produces a gradient, fails numerically, or is
    structurally unable to represent the task (\cref{sec:optimization}).
  \item[Gradient accuracy.] Central finite differences through each
    solver's own forward pass serve as the reference. We
    evaluate directional derivatives along $K$ random perturbation
    vectors ($K{=}6$--$20$) and report cosine similarity and relative
    $\ell_2$ error for all problems. For the two Navier--Stokes
    problems (F2, F3) we additionally compute the full Jacobian and
    inspect its singular value spectrum.
  \item[Performance.] Primal and VJP wall-clock time, their ratio, and
    peak (V)RAM consumption, measured on each solver's intended hardware
    target across three to four problem sizes (grid resolutions or cell counts).
    The VJP time is the cost of one gradient from scratch, a forward plus an
    adjoint pass (\cref{app:effort}).
    All timings are averaged over 3 runs after 1 warmup iteration
    (to pre-populate JIT caches and stabilize memory allocation).
  \item[Forward accuracy.] Resolution sweep against a reference solver,
    adherence to physical laws where applicable (e.g., solenoidal
    velocity). Where an analytical solution exists (F2, F3),
    precision is additionally measured against it.
  \item[Optimization convergence.] We run the optimizers with each solver's
    gradients on each benchmark task for a fixed iteration budget
    (500 iterations for F2 and F3; 1\,000 for H; 2\,500 for S) and report convergence
    curves and final objective value. Success is defined as reaching a
    final objective within 1\% of the best solution achieved across all
    solvers within the budget.
\end{description}

\subsection{Included solver backends}
\label{sec:solvers}

The solver set (\cref{tab:solvers}) is chosen to maximize diversity
along three axes: \emph{AD strategy} (autodiff, tape-based discrete
adjoints, hand-derived/kernel adjoints), \emph{numerical method} (projection-based, lattice Boltzmann,
finite-volume, spectral, finite elements with both direct and iterative
linear solvers), and \emph{solver abstraction level}.
These range from self-contained \emph{applications} (OpenFOAM),
through \emph{frameworks} that handle discretization and AD (FEniCS,
PhiFlow, JAX-CFD), to \emph{kernel toolkits} (NVIDIA Warp) where the
Mosaic wrapper \emph{is} the solver. Abstraction level determines
available gradient strategies and per-setup effort
(\cref{app:solvers}).

\begin{table}[t]
  \centering\small
  \caption{Solver backends included in the current release.
    Domain IDs match \cref{tab:domains}.}
  \label{tab:solvers}
  {\setlength{\tabcolsep}{4pt}
  \begin{tabular}{@{}llllll c@{}}
    \toprule
    \textbf{Solver} & \textbf{\shortstack[l]{Language /\\AD framework}}
      & \textbf{Strategy}
      & \textbf{\shortstack[l]{Discret-\\ization}} & \textbf{Numerics}
      & \textbf{GPU} & \textbf{Domains} \\
    \midrule
    JAX-CFD~\cite{kochkov2021machine}
      & Py/JAX & Auto & FD & Projection, RK4 & $\checkmark$ & F2 \\
    PhiFlow~\cite{holl2024phiflow}
      & Py/JAX & Auto+An. & FD & Central FD, Euler & $\checkmark$ & F2, F3 \\
    INS.jl~\cite{agdestein2024ins}
      & Julia/Zygote & Auto & FD & Projection, RK4 & & F2, F3 \\
    XLB~\cite{ataei2024xlb}
      & Py/JAX & Auto & LBM & Streaming & $\checkmark$ & F2, F3 \\
    PICT~\cite{franz2025pict}
      & Py/PyTorch & Auto+An. & FV & PISO, BDF1 & $\checkmark$ & F2, F3 \\
    Warp-NS~\cite{macklin2024warp}
      & Py/Warp & Auto+An. & FD & Projection, Euler & $\checkmark$ & F2, F3 \\
    Exponax~\cite{koehler2024apebench}
      & Py/JAX & Auto & Spectral & ETDRK2 & $\checkmark$ & F3 \\
%   FEniCS-NS~\cite{farrell2013automated}
%     & Py & Sym.+Auto & FE & IPCS, Crank-Nicolson & & F3 \\
    OpenFOAM~\cite{weller1998tensorial}
      & C++ & --- & FV & PISO/PIMPLE, Euler & & F2, F3 \\
    \midrule\midrule
    FEniCS~\cite{farrell2013automated}
      & Py & Sym.+Auto & FE & Direct (MUMPS) & & H, S \\
    Firedrake~\cite{rathgeber2016firedrake}
      & Py & Sym.+Auto & FE & Direct (MUMPS) & & H, S \\
    JAX-FEM~\cite{xue2023jaxfem}
      & Py/JAX & Auto+An. & FE & Direct (SuperLU) & $\checkmark$ & H, S \\
    TopOpt.jl~\cite{huang2021topoptjl}
      & Julia & Analytic & FE & Direct (CHOLMOD) & & S \\
    torch-fem~\cite{meyer2024torchfem}
      & Py/PyTorch & Auto+An. & FE & CG+Jacobi & $\checkmark$ & H, S \\
    deal.II~\cite{bangerth2007dealii,arndt2023dealii}
      & C++ & --- & FE & Direct (S), CG+SSOR (H) & & H, S \\
    \bottomrule
  \end{tabular}}
  \par\vspace{4pt}
  \begin{minipage}{\linewidth}\footnotesize\raggedright
    \textit{Strategy:}
    Auto\,=\,framework AD through discrete ops;
    An.\,=\,hand-derived adjoint of the discrete system;
    Sym.\,=\,adjoint derived at variational level via UFL;
    combinations indicate mixed strategies.
    \textit{Discr.:}
    FD\,=\,finite differences;
    FV\,=\,finite volume;
    FE\,=\,finite elements;
    LBM\,=\,lattice Boltzmann method.
    \textit{Numerics:}
    PISO\,=\,pressure-implicit splitting of operators;
    PIMPLE\,=\,merged PISO-SIMPLE;
    BDF1\,=\,1st-order backward differentiation;
    ETDRK\,=\,exponential time-differencing Runge--Kutta;
    Streaming\,=\,exact LBM streaming step.
  \end{minipage}
\end{table}

\paragraph{Inclusion and exclusion criteria}
A solver is included if it provides a general-purpose reverse-mode VJP
at cost proportional to a single forward solve. Forward-only solvers
serve as reference baselines. Excluded solvers and reasons are documented
in \cref{app:solvers}.

% ===========================================================================
\section{Experiments}
\label{sec:experiments}

The following experiments serve a dual purpose: they validate that
Mosaic's evaluation protocol surfaces practically relevant differences
between solvers, and they provide an initial empirical characterization
that practitioners can use to guide solver selection.
All experiments ran on an Azure Standard\_NC24ads\_A100\_v4 instance
(one NVIDIA A100 80\,GB PCIe, 24-core AMD EPYC 7V13; full spec in \cref{app:effort}).
Results reflect default solver configurations at the time of the pinned
release; all configurations are published and community-improvable via
pull request. Reproduction instructions are in \cref{app:software}.

\subsection{Benchmark summary and implementation effort}
\label{sec:baselines}

\Cref{tab:merged} summarizes benchmark results and gradient implementation
effort across all solvers; individual metrics are unpacked in the
subsections below.
Per-domain physical-accuracy validation, including the F2 cylinder forward
error and the divergence-RMS, kinetic-energy, and analytic-error sweeps for
the NS domains, is reported in \cref{app:physical_accuracy}.
Structural incompatibilities exclude some solvers from specific tasks;
causes are detailed in \cref{sec:optimization,app:solver-issues}.

Gradient implementation effort varies substantially among differentiable solvers.
We report bespoke lines of code (LOC) as a coarse proxy for integration
complexity (an imperfect metric, but one that captures order-of-magnitude
differences reproducibly).
Warp-NS requires over 800 bespoke lines for the initial-condition VJP alone,
a direct consequence of its low-level CUDA kernel model where each adjoint pass
must be written explicitly.
PhiFlow and JAX-CFD each require just 12 shared lines with no per-variable
additions, because autodifferentiation through their JAX backends provides
gradients at no extra implementation cost.
This pattern holds broadly: solvers built on modern AD frameworks (JAX,
PyTorch) consistently require zero per-variable extra lines, while hand-written
adjoint implementations carry substantially higher overhead.
The boundary is not always sharp: INS.jl uses Zygote reverse-mode AD yet
still requires over 470 shared lines of gradient-specific code,
illustrating that AD framework support does not always translate into
turnkey gradients.


\begin{table}[!htbp]
  \centering\small
  \setlength{\tabcolsep}{3pt}
  \caption{Benchmark results, gradient support, and integration effort
    per solver. Wall-clock times are means of 3 runs after 1 warmup.
    Full specifications in \cref{app:effort}; initial conditions in
    \cref{app:ics}. Legend below table.}
  \label{tab:merged}
  \vspace{4pt}
  \begin{tabular}{@{}l r r c c c c | ccc | r r r@{}}
    \toprule
    \textbf{Solver}
      & \multicolumn{2}{c}{\textbf{Wall time}}
      & \multicolumn{3}{c}{\textbf{Error}}
      & \textbf{Conv.}
      & \multicolumn{3}{c|}{\textbf{Grad.\ support}}
      & \multicolumn{3}{c@{}}{\textbf{Integ.\ effort}} \\
    \cmidrule(lr){2-3}\cmidrule(lr){4-6}\cmidrule(lr){8-10}\cmidrule(lr){11-13}
    & \textit{Fwd} & \textit{VJP} & \textit{Ref.} & \textit{An.} & \textit{FD} &
      & \multicolumn{3}{c|}{\textit{input variables}} & \textit{IO} & \textit{Num.} & \textit{Conf.} \\
    \midrule
    \multicolumn{7}{@{}l}{\textit{Navier--Stokes grid (F2), wall time in ms}}
      & \textit{$v_0$} & \textit{$\nu$} & \textit{inflow} & & & \\
    \cmidrule(lr){8-10}
    Warp-NS~\cite{macklin2024warp}        &    428 &    1400 & \textbf{1.9e-3} & \textbf{6.2e-4} & 8.4e-4 & $?$          & $\circ$   & $\circ$   & $?$     &   \textbf{0} & 1631 &  \textbf{0} \\
    INS.jl~\cite{agdestein2024ins}       &    343 &    2800 & 2.8e-3 & 2.4e-3 & \textbf{1.7e-5} & $\bot$       & $\bullet$ & $\times$  & $\bot$  &   117 & 1097 &  97 \\
    XLB~\cite{ataei2024xlb}               &     19 &    5200 & 2.5e-2 & 2.7e-2 & 2.1e-5 & $\checkmark$ & $\bullet$ & $\bullet$ & $\bullet$ &  13 &  464 &  58 \\
    PICT~\cite{franz2025pict}             &   2200 &    7200 & 2.4e-3 & 1.1e-3 & 1.2e-3 & $\checkmark$ & $\bullet$ & $\bot$    & $\bullet$ &  48 &  329 & 892 \\
    JAX-CFD~\cite{kochkov2021machine}     &     \textbf{11} &     \textbf{108} & 4.4e-3 & 4.6e-3 & 8.3e-4 & $\bot$       & $\bullet$ & $\bullet$ & $\bot$  &    38 &  286 &  21 \\
    PhiFlow~\cite{holl2024phiflow}        &    713 &    2200 & 2.8e-3 & 2.4e-3 & 2.2e-5 & $\triangle$  & $\bullet$ & $\bullet$ & $\bullet$ &  38 &  262 &  36 \\
    OpenFOAM~\cite{weller1998tensorial}   &   1300 &     $-$ & $-$    & 2.4e-3 & $-$    & $?$          & $?$       & $?$       & $?$     &   275 &  \textbf{127} & 437 \\
%   FEniCS~\cite{farrell2013automated}    &  10000 & 1248000 & 2.3e-3 & 1.5e-4 & 9.8e-4 & $\times$     & $\circ$   & $?$       & $\circ$ &   $-$ &  397 &  49 \\
    \midrule
    \multicolumn{7}{@{}l}{\textit{Navier--Stokes grid (F3), wall time in ms}}
      & \textit{$v_0$} & \textit{$\nu$} & & & & \\
    \cmidrule(lr){8-9}
    Warp-NS~\cite{macklin2024warp}        &    486 &    1500 & 2.0e-2 & \textbf{2.5e-3} & 2.5e-5 & $\checkmark$ & $\circ$   & $\circ$   & &   \textbf{0} & 1631 &  \textbf{0} \\
    INS.jl~\cite{agdestein2024ins}       &   2500 &   17000 & 4.1e-2 & 1.0e-2 & 9.9e-6 & $\checkmark$ & $\bullet$ & $\times$  & &   117 & 1097 &  97 \\
    XLB~\cite{ataei2024xlb}               &     43 &   48000 & 5.3e-2 & 1.9e-2 & \textbf{8.0e-6} & $\checkmark$ & $\bullet$ & $\bullet$ & &    13 &  464 &  58 \\
    PICT~\cite{franz2025pict}             &    767 &    3300 & \textbf{1.9e-2} & \textbf{2.5e-3} & 9.2e-5 & $\checkmark$ & $\bullet$ & $\bot$    & &    48 &  329 & 892 \\
    PhiFlow~\cite{holl2024phiflow}        &    627 &    1900 & 4.0e-2 & 1.0e-2 & 1.8e-5 & $\checkmark$ & $\bullet$ & $\bullet$ & &    38 &  262 &  36 \\
    Exponax~\cite{koehler2024apebench}    &     \textbf{22} &     \textbf{177} & 2.1e-2 & \textbf{2.5e-3} & 6.7e-4 & $\checkmark$ & $\bullet$ & $\bullet$ & &    41 &  160 &  \textbf{0} \\
    OpenFOAM~\cite{weller1998tensorial}   &   4400 &     $-$ & $-$    & \textbf{2.5e-3} & $-$    & $-$          & $?$       & $?$       & &   275 &  \textbf{127} & 437 \\
    \midrule
    \multicolumn{7}{@{}l}{\textit{Structural mechanics (S), wall time in s}}
      & \textit{$\rho$} & & & & & \\
    \cmidrule(l){8-8}
    deal.II~\cite{bangerth2007dealii}       &    166 &     $-$ & $-$    & $-$ & $-$    & $?$          & $\bot$    & & &   408 &  344 & 147 \\
    TopOpt.jl~\cite{huang2021topoptjl}     &     78 &  \textbf{15} & 0.0    & $-$ & 3.0e-5 & $\checkmark$ & $\circ$   & & &   124 &  155 &  \textbf{42} \\
    Firedrake~\cite{rathgeber2016firedrake} &    \textbf{7} &      17 & 0.0    & $-$ & 3.0e-5 & $\checkmark$ & $\circ$   & & &    39 &  221 & 137 \\
    FEniCS~\cite{farrell2013automated}     &    140 &      86 & 0.0    & $-$ & 3.0e-5 & $\checkmark$ & $\circ$   & & &   \textbf{0} &  222 &  63 \\
    JAX-FEM~\cite{xue2023jaxfem}           &     28 &      65 & 0.0    & $-$ & 3.0e-5 & $\checkmark$ & $\bullet$ & & &   \textbf{0} &  \textbf{150} & 133 \\
    \midrule
    \multicolumn{7}{@{}l}{\textit{Heat transfer (H), wall time in s}}
      & \textit{$\rho$} & \textit{src} & & & & \\
    \cmidrule(l){8-9}
    deal.II~\cite{bangerth2007dealii}       &    9.4 &          $-$ & $-$    & $-$ & $-$    & $?$          & $\bot$    & $\bot$    & &   381 &  263 &  75 \\
    FEniCS~\cite{farrell2013automated}      &     60 &     47 & 0.0    & $-$ & 2.2e-5 & $\checkmark$ & $\circ$   & $\circ$   & &   \textbf{0} &  513 &  \textbf{61} \\
    torch-fem~\cite{meyer2024torchfem}      &     33 &     11 & 0.0    & $-$ & 2.2e-5 & $\checkmark$ & $\bullet$ & $\bullet$ & &   \textbf{0} &  268 & 122 \\
    JAX-FEM~\cite{xue2023jaxfem}            &   \textbf{4.5} &    9.8 & 0.0    & $-$ & 2.2e-5 & $\checkmark$ & $\bullet$ & $\bullet$ & &   \textbf{0} &  \textbf{211} &  95 \\
    Firedrake~\cite{rathgeber2016firedrake} &     10 &     20 & 0.0    & $-$ & 2.2e-5 & $\checkmark$ & $\circ$   & $\circ$   & &   \textbf{0} &  420 & 102 \\
    \bottomrule
  \end{tabular}
  \par\vspace{4pt}
  \begin{minipage}{\linewidth}\footnotesize\raggedright
    \textit{Error:}
    Ref.\,=\,vs.\ reference solver;
    An.\,=\,vs.\ analytical solution (F2/F3);
    FD\,=\,vs.\ finite differences;
    $-$\,=\,not applicable.
    \textit{Conv.:}
    $\checkmark$\,=\,converged;
    $\triangle$\,=\,partial;
    $\times$\,=\,failed;
    $?$\,=\,not attempted;
    $\bot$\,=\,structurally excluded (\cref{app:solver-issues}).
    \textit{Grad.\ support:}
    $\bullet$\,=\,native AD;
    $\circ$\,=\,implemented by wrapper;
    $\times$\,=\,incorrect;
    $?$\,=\,not obtained;
    $\bot$\,=\,structurally blocked.
    \textit{Integ.\ effort} (bespoke LOC):
    IO\,=\,input/output marshalling;
    Num.\,=\,numerics \& gradient code;
    Conf.\,=\,solver configuration.
    PhiFlow partially converges on F2 ($\triangle$): drag is reduced but the optimizer stalls before reaching the XLB/PICT solution.
  \end{minipage}
\end{table}

\subsection{Gradient quality}
\label{sec:quality}

To assess whether AD gradients are correct, we compare each solver's
gradients against central finite differences at solver-specific optimal
perturbation size.

\begin{figure}[t]
  \centering
  \includegraphics[width=\linewidth]{figures/horizon_sweep_limits.pdf}
  \caption{VJP cost and stability vs.\ rollout length on F3.}
  \label{fig:horizon_sweep}
\end{figure}

All differentiable solvers reach cosine similarity $>0.999$ between AD
and central finite-difference gradients at their optimal perturbation
size~$\varepsilon$, confirming AD correctness; per-solver minimum FD
errors are reported in the \textit{Error (FD)} column of
\cref{tab:merged}. The optimal $\varepsilon$ shifts by orders of
magnitude across solvers and with rollout length and viscosity; see
\cref{app:gradient} for the full per-domain U-shape curves and
discussion.

While gradient accuracy is uniformly high, \emph{conditioning} is not.
\Cref{fig:conditioning} reports Jacobian conditioning via normalized SVD
spectra on the 3D periodic NS problem. Conditioning varies dramatically
across solvers at the same resolution, reflecting the numerical method
rather than the AD strategy: the same governing equations discretized
with lattice Boltzmann versus pressure projection yield radically
different spectra.
The sharp spectral drop-off in projection-based solvers is attributable
to singular vectors associated with non-solenoidal velocity modes, which
incompressible flow forbids by construction~\citep{guermond2006overview,bhatia2013helmholtz}.
This distinction matters for optimization: conditioning predicts which
solvers will struggle with second-order methods and long-horizon
problems (\cref{sec:scaling,sec:optimization}).
\begin{figure}[t]
  \centering
  \includegraphics[width=\linewidth]{figures/recovery_main.pdf}
  \caption{F3 IC recovery for selected solvers and different optimizers; full version in \cref{fig:app:recovery}.}
  \label{fig:fd_verification}
\end{figure}

\begin{figure}[t]
  \centering
  \includegraphics[width=\linewidth]{figures/jacobian_svd_comparison.pdf}
  \caption{Normalized Jacobian singular value spectra ($\sigma_i / \sigma_0$) for six solvers on the 3D periodic NS problem at four rollout lengths.
    Exponax and XLB maintain a relatively flat spectrum across all rollout lengths; INS.jl and PhiFlow develop very steep spectra, indicating large condition numbers $\kappa = \sigma_1/\sigma_n$.}
  \label{fig:conditioning}
\end{figure}

\subsection{Computational cost and scaling}
\label{sec:scaling}

Absolute solver cost varies by orders of magnitude at the same
resolution (\cref{tab:merged}; full scaling curves in \cref{app:cost}),
but speed and generality are entangled: the fastest backends
(Exponax, JAX-CFD) support only periodic boundaries and no solid
obstacles, while solvers that handle the full F2 drag task (XLB, PICT)
are up to ${\sim}200\times$ slower. Cost comparisons are therefore meaningful only
within the set of solvers that can represent a given task (\cref{sec:optimization}).

At larger resolutions and longer rollouts, memory becomes the binding
constraint before gradient accuracy does: several solvers exceed the
16\,GB GPU memory limit and fail to produce gradients entirely
(\cref{fig:horizon_sweep}).
Among solvers that remain within budget, the cost of obtaining a
gradient depends on the AD strategy.
Source-transformation AD (JAX-based solvers) maintains a nearly constant
VJP/forward ratio across problem sizes, while tape-based adjoints
(FEniCS, Firedrake) scale more steeply (\cref{app:cost}), making AD
strategy a first-order predictor of gradient cost at scale.
All timings are measured on one A100 GPU;
relative rankings may differ on other hardware, particularly between
GPU and CPU solvers.

\subsection{Optimization task performance}
\label{sec:optimization}

All solvers share the same hyperparameters and iteration budget per
task (settings in \cref{app:recovery}), so differences in convergence
are attributable to the solver's gradient alone.
On S, all solvers converge to the same optimum. On H, every solver
drives the temperature misfit down by orders of magnitude, but the poorly
conditioned inversion leaves the recovered conductivity dependent on the
optimizer's path within the iteration budget
(\cref{fig:app:conductivity_recovery}).
On F3 (IC recovery), we report the relative $\ell_2$ distance between
the recovered and true initial states (\textit{IC error})
alongside the optimization loss. Adam stagnates
above 40\% IC error across all solvers (\cref{fig:fd_verification}):
the loss decreases, but the recovered IC differs noticeably from the
target.
L-BFGS converges to substantially lower loss and below 6\% IC error in
fewer iterations.
Spectrally projecting the gradient onto the divergence-free subspace
before each L-BFGS update (L-BFGS+proj) drops the IC error below 0.5\%,
because the projection prevents non-solenoidal components from
corrupting the curvature model.
The 2D drag minimization (F2) is the hardest deployment test, and
also the task where structural incompatibilities thin the solver field
most: only three of seven backends support the non-periodic, obstacle-laden
setup required for drag optimization.
Of those, XLB and PICT converge to consistent drag reductions of
around 60\%; PhiFlow partially converges, reducing drag by ${\sim}35\%$
before stalling. The remaining backends are excluded for structural
reasons (periodic BCs, penalization incompatibilities;
\cref{app:solver-issues}). This attrition itself is a key finding:
on the most practically demanding task, solver selection is constrained
by setup compatibility before gradient quality enters the picture.

% ===========================================================================
\section{Discussion and conclusion}
\label{sec:conclusion}

Gradient accuracy, computational cost, Jacobian conditioning, and setup
compatibility all vary across solvers, across domains, and across
operating regimes within a single domain. Yet these differences do not
affect optimization outcomes uniformly. Based on the four benchmark
tasks evaluated here, we highlight five practitioner-facing observations.

\paragraph{Gradient quality is often good enough; the real bottlenecks
lie elsewhere}
On every task where they can run, all solvers that produce gradients
converge to equivalent optima despite order-of-magnitude differences in
accuracy and conditioning. While more complex tasks may reveal cases
where gradient quality becomes limiting, the pattern on these benchmarks
suggests that conditioning is likely a better predictor of
training stability than pointwise gradient accuracy.

\paragraph{Rollout length is the critical scaling parameter}
Jacobian conditioning degrades with rollout length across all solvers
(\cref{fig:conditioning}), with projection-based solvers deteriorating
faster than LBM or spectral methods. Memory scales likewise, and
several solvers exceed GPU limits before gradient quality becomes
relevant.

\paragraph{On the negative-control domains, pick the fastest solver}
On the negative-control domains (H, S), all solvers produce equivalent
results regardless of AD strategy or numerical method. Because gradient
quality is not a differentiator in these regimes, the selection criterion
reduces to forward-solve speed and integration effort.

\paragraph{Reverse-mode AD offers the best effort-to-capability
tradeoff}
Hand-derived adjoints can be faster than the forward solve
(\cref{tab:merged}), but each new control variable or objective requires
additional implementation work. Reverse-mode AD through modern
frameworks (JAX, PyTorch) produces gradients accurate enough for
optimization at a fraction of the implementation cost.

\paragraph{What Mosaic enables beyond this evaluation}
Mosaic's uniform interface makes each backend a drop-in differentiable
simulator for solver-in-the-loop or hybrid physics-ML workflows,
and the template mechanism lets the benchmark grow with the community.

\paragraph{Limitations}
The evaluation rests on three assumptions: that central
finite-difference gradients at optimal perturbation size are a reliable
reference (\cref{fig:app:fdcheck}), that out-of-the-box solver configurations
are representative of practitioner experience, and that standalone
optimization tasks are a meaningful proxy for gradient quality in broader
pipelines. Gradient quality may differ on untested setups; the
template mechanism is the mitigation. We benchmark solver gradients in
isolation, without learned components, so hybrid pipelines may introduce
failure modes not captured here. Results are pinned to specific solver
versions; upstream updates may change gradient behavior.

\paragraph{Future work}
Natural extensions include additional domains (multiphysics, compressible
flow), solver backends (Enzyme/LLVM-level AD), and hybrid physics-ML tasks.
The infrastructure is in place, and we invite community contributions.

% ===========================================================================

% ── REFERENCES ──────────────────────────────────────────────
\PLrefheading
\bibliography{refs}

\clearpage

% ===========================================================================
\appendix

\section{Software and reproducibility}
\label{app:software}

\subsection{License, access, and reproduction}

Mosaic is released under the Apache~2.0 license (individual solver backends
retain their upstream licenses). The repository, documentation, and
reproduction instructions are available at
\url{https://github.com/pasteurlabs/mosaic}.

The paper corresponds to the tag
\href{https://github.com/pasteurlabs/mosaic/tree/v0.1%2Bpaper-repro}{\texttt{v0.1+paper-repro}}
which contains all code, plotting scripts,
solver wrappers, and a pinned dependency lockfile sufficient for full
reproduction. Two reproduction paths are supported. Figures can be
regenerated from the result archive linked in the repository README without
re-running any solver. Alternatively, the full experiment suite can be
re-run from scratch; this requires Docker and, for GPU solvers, an NVIDIA
GPU with the NVIDIA Container Toolkit. The lockfile pins every transitive
dependency to the exact versions used for the paper, and solver containers
are built on-demand by the harness. Step-by-step instructions for both
paths are in the repository README.

\subsection{Architecture}

Each solver backend is a \emph{Tesseract}~\cite{haefner2025tesseract}: a
Docker container that exposes a typed forward map (\texttt{apply}), shape
inference (\texttt{abstract\_eval}), and optionally a VJP
(\texttt{vector\_jacobian\_product}) over Pydantic schemas. The container
boundary isolates solver-specific dependencies (JAX, PyTorch, Julia, C++,
Warp) from one another and from the evaluation harness.
\texttt{tesseract-jax} bridges the container interface into JAX's
\texttt{jit}/\texttt{grad} transform chain, so the harness can
differentiate through any backend regardless of its native AD framework
using a single code path.

Each backend is defined by three files: a Python module implementing the
entry points above, a YAML configuration declaring solver metadata (name,
AD strategy, backend framework, plot style), and a requirements file listing
dependencies. The evaluation harness discovers solvers automatically from
their configuration metadata and runs five benchmark suites (forward
accuracy, gradient quality, computational cost, optimization convergence,
and initial-condition visualization) without solver-specific code.

\subsection{Extensibility}

Contributing a new solver requires implementing the three files described
above for an existing domain. The harness discovers the backend
automatically from its configuration metadata; no registration step or
evaluation-code changes are needed. A step-by-step tutorial (adding an LBM
solver to the fluid domain) is included in the repository documentation.
New benchmark domains can be scaffolded from built-in templates and require
interface schemas, at least one solver backend, and a problem configuration
defining initial conditions, error metrics, and suite defaults. Existing
solver configurations can also be improved via pull requests; CI re-runs
the full evaluation suite and posts a before/after comparison.

\subsection{Maintenance and versioning}

We maintain the evaluation protocol, the Tesseract interface contract, CI
infrastructure, and the paper-pinned solver set (tagged release, kept
buildable). New solver contributions, domain-specific task formulations,
and upstream solver maintenance are community-owned. If an upstream solver
breaks due to a dependency update, we document version pins and accept
community fixes but do not patch upstream code ourselves.

Each release is tagged with a semantic version and accompanied by a pinned
dependency lockfile. Continuous integration runs the benchmark suite on
every push to
\texttt{main} that touches a solver or the harness, committing results
to a dedicated branch with a machine-readable diff against the previous
baseline. Tagged releases run the full suite and publish a result archive
as a GitHub release artifact. This keeps published results current as
solvers evolve and new backends land.

\section{Extended experimental results}
\label{app:extended}

\subsection{Benchmark specifications}
\label{app:effort}

\paragraph{Integration effort}
The \textbf{IO}, \textbf{Num.}, and \textbf{Conf.}\ columns in \cref{tab:merged}
report bespoke lines of code by category.
\emph{IO} covers serialization and file-based data exchange with external solver processes.
\emph{Num.}\ is problem-specific numerical code (assembly, time-stepping, boundary
conditions, and gradient plumbing).
\emph{Conf.}\ covers solver initialization and mesh setup.
Glue code (Tesseract schema definitions and entry-point boilerplate) is excluded; it
consistently amounts to 100--400\,LOC across all backends.

\paragraph{Error column definitions}
\textit{Ref.}: $\|\mathbf{f}_{\text{solver}} - \mathbf{f}_{\text{ref}}\|_2 / \|\mathbf{f}_{\text{ref}}\|_2$,
where $\mathbf{f}_{\text{ref}}$ is the forward state from the reference solver (OpenFOAM for F2/F3; deal.II for S/H).
This is a forward accuracy metric, not a gradient comparison.
\textit{An.}: $\|\mathbf{f}_{\text{solver}} - \mathbf{f}_{\text{TGV}}\|_2 / \|\mathbf{f}_{\text{TGV}}\|_2$,
where $\mathbf{f}_{\text{TGV}}$ is the TGV closed-form forward solution (F2/F3 only).
\textit{FD}: minimum over $\varepsilon \in \{10^{-6}, \ldots, 10^{-1}\}$ of
$\|\mathbf{g}_{\text{solver}} - \mathbf{g}_{\text{FD},\varepsilon}\|_2 / \|\mathbf{g}_{\text{FD},\varepsilon}\|_2$,
where $\mathbf{g}_{\text{FD},\varepsilon}$ is a central finite-difference gradient estimate at step size $\varepsilon$.

\paragraph{Hardware and timing protocol}
All experiments ran on an Azure Standard\_NC24ads\_A100\_v4 instance
(24 vCPUs of an AMD EPYC 7V13, 216\,GiB RAM, one NVIDIA A100 80\,GB PCIe);
CPU-only solvers ran on the same instance's CPU.
The results on the project website are regenerated by continuous integration
on NVIDIA T4 runners, so their absolute times differ from the ones reported here.
Wall times use one warmup call (excluded, absorbs JIT compilation and CUDA kernel
caching) followed by three timed trials; the reported value is the mean.
A forward time covers one \texttt{apply} request and a VJP time one
\texttt{vector\_jacobian\_product} request, whose cotangent comes from an untimed
\texttt{apply} beforehand. Every solver recomputes its forward solution inside the
VJP request, so the VJP time is the cost of one gradient from scratch: a forward
plus an adjoint pass. TopOpt.jl's compliance gradient is self-adjoint and needs no
second solve, so its VJP costs about one forward solve.
The timed window covers host-side serialization, the Tesseract HTTP round-trip, solver
compute, and response deserialization; container startup is excluded.

\paragraph{F2 (2D NS, grid)}
Periodic domain $[0,2\pi)^2$, uniform Cartesian grid, Taylor--Green vortex IC\@.
\textbf{Fwd/VJP}: $64\times64$ grid (4\,096 cells), $\nu=0.01$, $\Delta t=0.01$, 100 steps.
\textbf{Error (Ref., An.)}: agreement experiment, $64\times64$ grid, $\nu=0.01$,
$\Delta t=0.05$, 20 steps; Ref.\ vs.\ fine JAX-CFD reference ($\Delta t=0.01$, 100 steps),
An.\ vs.\ TGV closed-form $\mathbf{u}(t)=\mathbf{u}(0)e^{-2\nu t}$.
\textbf{Error (FD)}: finite-difference gradient check, $16\times16$ grid, multimode IC,
$\nu=0.001$, $\Delta t=0.05$, 20 steps, 20 random directions.
\textbf{Convergence}: drag-reduction optimization, $32\times32$ grid, channel domain
$[0,1]^2$, cylinder at $(0.5,0.5)$ radius $0.05$, Re\,=\,20 ($\nu=0.0025$),
$\Delta t=0.02$, 400 steps, Adam ($\eta=5\times10^{-4}$), up to 500 iterations.
An L-BFGS variant with the same physics is reported in \cref{app:recovery}.

\paragraph{F3 (3D NS, grid)}
Periodic domain $[0,2\pi)^3$, uniform Cartesian grid, TGV IC\@.
\textbf{Fwd/VJP}: $16^3$ grid (4\,096 cells), $\nu=0.01$, $\Delta t=0.01$, 50 steps.
\textbf{Error (Ref., An.)}: forward agreement, $16^3$ grid, $\nu=0.01$, $\Delta t=0.01$,
50 steps; Ref.\ vs.\ fine Exponax reference ($\Delta t=0.002$, 250 steps),
An.\ vs.\ TGV closed-form.
\textbf{Error (FD)}: finite-difference gradient check, $\nu=0.001$, $\Delta t=0.05$,
10 steps, 10 random directions.
\textbf{Convergence}: IC recovery, $\nu=0.01$, $\Delta t=0.02$, 100 steps,
Adam ($\eta=10^{-3}$), up to 500 iterations.

\paragraph{S (Structural mechanics, mesh)}
Cantilever beam $[0,2]\times[0,1]\times[0,1]$, HEX8 elements, SIMP penalization
($p=3$, $E_\text{max}=70\,000$\,MPa, $\nu_\text{mat}=0.3$, $x_\text{min}=10^{-3}$).
\textbf{Fwd/VJP}: $128\times2\times64$ mesh (${\approx}16\text{k}$ elements), uniform
density $\rho_0=0.5$, distributed unit load $F=1$.
\textbf{Error (FD)}: $8\times2\times4$ mesh (64 elements), corner point load,
6 random directions.
\textbf{Convergence}: compliance minimization, $16\times2\times8$ mesh (256 elements),
volume fraction $0.5$, Adam ($\eta=5\times10^{-2}$), up to 2\,500 iterations.

\paragraph{H (Heat transfer, mesh)}
Quasi-2D heated slab $[0,2]\times[0,1]$, HEX8 elements (single layer), SIMP
penalization ($p=3$, $k_\text{max}=1$, $k_\text{min}/k_\text{max}=10^{-3}$).
\textbf{Fwd/VJP}: $256\times128\times1$ mesh (${\approx}33\text{k}$ elements), uniform
conductivity $k_0=0.5$, unit distributed heat source.
\textbf{Error (FD)}: $8\times4\times1$ mesh (32 elements), 6 random directions.
\textbf{Convergence}: conductivity recovery from two-Gaussian target,
$16\times8\times1$ mesh, L-BFGS, up to 1\,000 iterations.

\subsection{Initial conditions}
\label{app:ics}

Each benchmark domain uses a fixed set of initial conditions, shared across
all solvers.

\paragraph{Incompressible NS (\cref{fig:app:ics_ns})}
\textbf{Taylor--Green vortex (TGV, 2D):}
$u = \sin(x)\cos(y)$, $v = -\cos(x)\sin(y)$ on $[0,2\pi)^2$.
Admits the analytic solution $\mathbf{u}(t)=\mathbf{u}(0)e^{-2\nu t}$,
enabling exact error measurement without a reference solver.
\textbf{Multimode:} solenoidal field constructed from a spectral-space
stream function with energy concentrated in a ring at wavenumber $k=2$
($\sigma_k=0.5$) and random phases; peak speed normalized to~0.3.
\textbf{Uniform inflow} and \textbf{flat inlet profile} are spatially
constant fields; not shown.
\textbf{Taylor--Green vortex (TGV, 3D):}
$u = \sin(x)\cos(y)\cos(z)$, $v = -\cos(x)\sin(y)\cos(z)$, $w = 0$
on $[0,2\pi)^3$.
Evolves into a vortex-dominated turbulent state under the NS dynamics.
\textbf{Arnold--Beltrami--Childress (ABC):}
$u = A\sin(z)+C\cos(y)$, $v = B\sin(x)+A\cos(z)$,
$w = C\sin(y)+B\cos(x)$ with $A=B=C=1$ on $[0,2\pi)^3$.
A steady Euler solution with chaotic particle trajectories.

\begin{figure}[!htbp]
  \centering
  \includegraphics[width=\linewidth]{figures/appendix_ics_ns_combined.pdf}
  \caption{Initial conditions for the F2 and F3 domains (vorticity; 3D shown as $z{=}0$ slice).}
  \label{fig:app:ics_ns}
\end{figure}

\paragraph{Structural mechanics (\cref{fig:app:ics_structural})}
All ICs are material density fields over the beam mesh.
\textbf{Uniform:} $\rho_0 = 0.5$ everywhere.
\textbf{Random:} Gaussian noise centered at $\rho_0=0.5$ ($\sigma=0.3$),
clipped to $[0.05,\,0.95]$.
\textbf{Two density bumps:} two Gaussian pillars ($\rho_\text{peak}=0.95$,
$\sigma=0.12\,L$) at $x=0.35\,L_x$ and $x=0.75\,L_x$, on a soft
background ($\rho_\text{bg}=0.1$).

\begin{figure}[!htbp]
  \centering
  \includegraphics[width=\linewidth]{figures/appendix_ics_structural_mesh.pdf}
  \caption{Initial conditions for the structural mechanics domain.}
  \label{fig:app:ics_structural}
\end{figure}

\paragraph{Heat transfer (\cref{fig:app:ics_thermal})}
All ICs are conductivity or source fields over the thermal mesh.
\textbf{Uniform conductivity:} $k_0 = 0.5$ everywhere.
\textbf{Random conductivity:} Gaussian noise centered at $k_0=0.5$
($\sigma=0.3$), clipped to $[0.05,\,0.95]$.
\textbf{Gaussian source:} single heat source centered at
$(0.5\,L_x,\,0.5\,L_y)$ with width $\sigma=0.2\,\min(L_x,L_y)$.
\textbf{Two-Gaussian conductivity:} ground-truth field for the conductivity
recovery task; two Gaussian peaks at $(0.3\,L_x,\,0.5\,L_y)$ and
$(0.7\,L_x,\,0.5\,L_y)$ with width $\sigma=0.15\,\min(L_x,L_y)$.

\begin{figure}[!htbp]
  \centering
  \includegraphics[width=\linewidth]{figures/appendix_ics_thermal_mesh.pdf}
  \caption{Initial conditions for the heat transfer domain.}
  \label{fig:app:ics_thermal}
\end{figure}

\subsection{Cost}
\label{app:cost}

\Cref{fig:app:scaling,fig:app:scaling:fem} show forward time, VJP time, and VJP/forward ratio vs.\ element count for all four domains.
The x-axis is total element count, parameterized by a single integer $N$:
F2 uses an $N\times N$ uniform Cartesian grid ($N\in\{64,128,192,256\}$, giving $4\text{k}$--$66\text{k}$ cells);
F3 uses an $N\times N\times N$ Cartesian grid ($N\in\{16,32,48,64\}$, giving $4\text{k}$--$262\text{k}$ cells);
S uses an $N\times 2\times\lfloor N/2\rfloor$ hexahedral mesh ($N\in\{8,\ldots,3200\}$, giving $64$--$10\text{M}$ elements);
H uses an $N\times\lfloor N/2\rfloor\times 1$ hex mesh ($N\in\{16,\ldots,4500\}$, giving $128$--$10\text{M}$ elements).
Table~\ref{tab:merged} wall times are at $N=64$ (F2), $N=16$ (F3), $N=128$ (S), $N=256$ (H).

The fixed RPC overhead (host serialization and HTTP round-trip) adds approximately 2--20\,ms per call and is visible as a floor for the fastest solvers at small problem sizes.
Lines that terminate before the full size range indicate that the solver either raised an exception (typically an out-of-memory error) or exceeded the 1\,000\,s per-trial wall limit at that problem size (the underlying Tesseract HTTP request has a separate 1\,200\,s watchdog); all larger sizes for that solver are omitted.
Forward times span several orders of magnitude across solver families at identical resolution, reflecting differences in algorithmic complexity and hardware target (CPU vs.\ GPU).
VJP overhead tracks forward cost for source-transformation AD (JAX).
For the NS domains the VJP/forward ratio is roughly constant across problem sizes: both forward and backward passes integrate the same $T$ timesteps, so their costs scale identically with $N$.
For structural and thermal domains the ratio sits close to~1 at moderate sizes because the adjoint of a symmetric elliptic system shares the same stiffness matrix as the forward solve and can reuse its factorization; at the largest problem sizes the adjoint-based solvers (FEniCS, Firedrake) show steeper VJP scaling as the adjoint solve cost begins to dominate.

\begin{figure}[!htbp]
  \centering
  \includegraphics[width=\linewidth]{figures/scaling_ns_combined.pdf}
  \caption{Gradient overhead scaling for the fluid domains (2D NS, top;
    3D NS, bottom). Left: forward time; center: VJP time; right:
    VJP/forward ratio (dashed line marks 1:1). All axes are log-log.
    (G)\,/\,(C) denotes GPU\,/\,CPU execution.}
  \label{fig:app:scaling}
\end{figure}

\begin{figure}[!htbp]
  \centering
  \includegraphics[width=\linewidth]{figures/scaling_fem_combined.pdf}
  \caption{Gradient overhead scaling for the mesh domains (Structural,
    top; Thermal, bottom). Same layout as \cref{fig:app:scaling}.
    deal.II VJP times are omitted (no native adjoint).
    Source-transformation AD (JAX-FEM) maintains a nearly constant
    ratio; adjoint-based solvers (FEniCS, Firedrake) show steeper VJP
    scaling at large problem sizes.}
  \label{fig:app:scaling:fem}
\end{figure}

\subsection{Physical accuracy}
\label{app:physical_accuracy}

As a basic sanity check, we run agreement and baseline experiments on each NS domain: agreement measures relative $\ell_2$ error vs.\ the analytic TGV solution across a $\nu$ sweep at fixed resolution, and baseline measures convergence with resolution at one timestep.
All solvers that support the periodic TGV domain pass both checks with errors below $10^{-2}$.

\paragraph{F2 cylinder-flow forward accuracy}
The F2 domain uses a cylinder obstacle in a channel.
Unlike the periodic TGV, no analytic solution exists, so solvers are compared against the cross-solver consensus (mean over valid solvers).
This experiment uses a unit-square domain $[0,1]^2$ with a cylinder at $(0.5, 0.5)$, radius $0.1$, $N=64$, 500 steps at $\Delta t=0.01$; the physical accuracy sweeps below use the full channel domain $[0,8]\times[0,2]$. The F2 drag-optimization convergence experiment (\cref{app:effort}) uses the same unit-square domain but a smaller cylinder (radius $0.05$, $D=0.1$), which sets $\text{Re}=U D/\nu=20$ at the fixed $U=0.5$, $\nu=0.0025$ used there; both setups are F2 cylinder-channel flows but exercise different aspects of the benchmark.
Three solvers (JAX-CFD, INS.jl, Warp-NS) produce invalid results on this geometry due to boundary condition incompatibilities.
\Cref{fig:app:cylinder} shows consensus error vs.\ $\nu$ and final vorticity fields for the four solvers that run successfully.

\begin{figure}[!htbp]
  \centering
  \includegraphics[width=\linewidth]{figures/appendix_cylinder.pdf}
  \caption{F2 cylinder-flow forward accuracy: relative $\ell_2$ error vs.\ the cross-solver consensus as a function of $\nu$ at $N=64$, 500 steps (left), and final vorticity fields at $\nu=0.01$ (right). Solvers that fail on this geometry are omitted.}
  \label{fig:app:cylinder}
\end{figure}

\paragraph{Physical accuracy sweeps}
For the two fluid domains, \cref{fig:app:pa_ns2d,fig:app:pa_ns3d} show resolution~$N$, viscosity~$\nu$, and rollout length against three physical metrics.
F2 uses the $[0,8]\times[0,2]$ cylinder-channel at Re\,=\,20, sweeping $N\in\{16,32,64,128\}$, $\nu\in\{0.001,0.005,0.01,0.05,0.1\}$, and up to 100 steps.
F3 uses the $[0,2\pi]^3$ TGV box at $N=16$, sweeping $\nu\in\{0.001,0.01,0.05,0.1\}$ and up to 50 steps.
\textbf{Analytic TGV error} is the relative $\ell_2$ error against the Taylor--Green vortex closed-form solution. Lower is better, with values below $10^{-2}$ indicating good accuracy.
\textbf{Divergence RMS} measures how strongly the solver violates the incompressibility constraint $\nabla \cdot \mathbf{u} = 0$. For a well-posed incompressible solver this should be near machine precision ($\lesssim 10^{-6}$), while values above $10^{-2}$ indicate a solver that does not enforce incompressibility exactly (e.g.\ LBM methods).
\textbf{Kinetic energy} is the domain-averaged $\frac{1}{2}\|\mathbf{u}\|^2$.
For 2D NS the analytical ground truth is $\mathrm{KE}(t) = \tfrac{1}{4}e^{-4\nu t}$, shown as a dashed reference line, and solvers should track it closely.
For 3D NS no closed-form solution exists, so all solvers should at least agree with each other; spread at a given parameter value indicates a physical accuracy problem.
Solvers that enforce the divergence-free constraint exactly at each step by construction (PhiFlow, Warp-NS, PICT, JAX-CFD, INS.jl, Exponax) sit at or near machine precision regardless of resolution or viscosity.
XLB exhibits a compressibility floor in divergence RMS and analytic error that is largely independent of $\nu$ or resolution, consistent with the $O(\mathrm{Ma}^2)$ compressibility error inherent to the LBM formulation.

For the mesh domains, all structural and thermal solvers agree to machine precision across the full load range, consistent with the expected $C \propto F^2$ scaling of linear compliance.

\begin{figure}[!htbp]
  \centering
  \includegraphics[width=\linewidth]{figures/appendix_pa_ns2d.pdf}
  \caption{2D NS physical accuracy. Each panel sweeps one parameter (resolution, $\nu$, or steps) on the x-axis against one physical metric on the y-axis.}
  \label{fig:app:pa_ns2d}
\end{figure}

\begin{figure}[!htbp]
  \centering
  \includegraphics[width=\linewidth]{figures/appendix_pa_ns3d.pdf}
  \caption{3D NS physical accuracy on the triply-periodic TGV domain $[0,2\pi)^3$.
    Rows sweep resolution $N$ (top), viscosity $\nu$ (middle), and rollout length (bottom).
    Columns report analytic TGV error, divergence RMS, and kinetic energy.
    No closed-form kinetic-energy reference exists for 3D NS, so solvers are compared
    against each other; spread at a given parameter value indicates a physical accuracy
    discrepancy.}
  \label{fig:app:pa_ns3d}
\end{figure}

\subsection{Gradient quality}
\label{app:gradient}

\paragraph{Finite-difference verification}
\Cref{fig:app:fdcheck} extends the FD verification from \cref{sec:quality} to all four domains.
The U-shaped error curve (truncation error at large~$\varepsilon$, floating-point roundoff at small~$\varepsilon$) is visible in every domain and for every solver that provides gradients.
The optimal~$\varepsilon$ shifts noticeably between solver families and between domains: spectral solvers (Exponax) tolerate larger perturbations than projection-based solvers, and structural/thermal solvers operate at a different optimal scale than fluid solvers.
All differentiable solvers reach cosine similarity $>0.999$ at their optimal~$\varepsilon$, confirming gradient correctness across domains.

\begin{figure}[!htbp]
  \centering
  \includegraphics[width=\linewidth]{figures/appendix_fd_check.pdf}
  \caption{FD gradient verification across all four domains.
    Top row: relative $\ell_2$ error vs.\ perturbation size~$\varepsilon$.
    Bottom row: subspace cosine similarity vs.~$\varepsilon$.
    All differentiable solvers achieve cosine similarity $>0.999$ at their optimal~$\varepsilon$.}
  \label{fig:app:fdcheck}
\end{figure}

\begin{figure}[!htbp]
  \centering
  \includegraphics[width=0.85\linewidth]{figures/fd_check.pdf}
  \caption{FD gradient verification on the 2D NS domain (focused view of the leftmost column of \cref{fig:app:fdcheck}).
    Relative $\ell_2$ error (left) and cosine similarity (right) vs.\ perturbation size~$\varepsilon$, normalized by the RMS magnitude of the initial condition.
    The U-shaped curve reflects the truncation/roundoff tradeoff; optimal~$\varepsilon$ varies by orders of magnitude across solvers.}
  \label{fig:app:fdcheck:ns2d}
\end{figure}

\paragraph{Jacobian singular value spectra}
\Cref{fig:app:svd2d} shows normalized singular value spectra $\sigma_i/\sigma_0$ of the solver Jacobian for 2D NS across varying viscosity and rollout length; the 3D NS spectra appear in \cref{fig:conditioning} in the main text.
Faster spectral decay indicates a more ill-conditioned Jacobian and correspondingly harder gradient-based optimization.
Across both domains, increasing rollout length ($T$) consistently steepens the decay, while higher viscosity ($\nu$) slightly moderates it.
Solver families cluster: spectral solvers retain flatter spectra at longer horizons, while projection-based solvers exhibit sharper decay.
The sharp spectral drop-off in projection-based solvers is attributable to singular vectors associated with non-solenoidal velocity modes, which incompressible flow forbids by construction~\citep{guermond2006overview,bhatia2013helmholtz}; these higher modes are mostly attributed to divergence and are correspondingly less relevant for gradient-based optimization.

\begin{figure}[!htbp]
  \centering
  \includegraphics[width=\linewidth]{figures/appendix_jacobian_svd_2d.pdf}
  \caption{Normalized Jacobian singular value spectra for 2D NS.
    Each panel shows a different $(\nu, T)$ configuration.
    Faster decay corresponds to a more ill-conditioned Jacobian and harder gradient optimization.}
  \label{fig:app:svd2d}
\end{figure}

\paragraph{Gradient quality over rollout length}
\Cref{fig:app:horizon} shows how FD gradient quality evolves with rollout length on the 2D NS domain.
Even solvers with correct gradients at short horizons can degrade at longer rollouts, consistent with the Jacobian conditioning analysis in \cref{fig:conditioning}: solvers with rapidly decaying singular value spectra lose useful gradient signal faster.
The optimal perturbation size~$\varepsilon$ also drifts with horizon, reinforcing that a fixed FD step used for validation can give misleading results at longer rollout lengths.

\begin{figure}[!htbp]
  \centering
  \includegraphics[width=\linewidth]{figures/horizon_sweep.pdf}
  \caption{Gradient quality vs.\ rollout length (2D NS).
    Left: gradient norm over rollout length.
    Middle: best-achievable FD relative error at each horizon.
    Right: optimal perturbation size~$\varepsilon$.
    Solvers with poor Jacobian conditioning degrade fastest with increasing rollout length.}
  \label{fig:app:horizon}
\end{figure}

\paragraph{FD U-curves}
\Cref{fig:app:ucurves_f2,fig:app:ucurves_f3} show full FD U-curves (relative error vs.\ perturbation size~$\varepsilon$) at each rollout length, for the 2D and 3D NS domains respectively.
Each U-curve has a characteristic minimum: too-small $\varepsilon$ is dominated by floating-point cancellation; too-large $\varepsilon$ by nonlinear truncation error.
The minimum shifts rightward and rises as the rollout grows, reflecting increasing Jacobian ill-conditioning.

\begin{figure}[!htbp]
  \centering
  \includegraphics[width=\linewidth]{figures/appendix_ucurves_f2.pdf}
  \caption{FD U-curves for F2 (2D NS): relative FD error vs.\ perturbation size~$\varepsilon$ at each rollout length~$T$.}
  \label{fig:app:ucurves_f2}
\end{figure}

\begin{figure}[!htbp]
  \centering
  \includegraphics[width=\linewidth]{figures/appendix_ucurves_f3.pdf}
  \caption{FD U-curves for F3 (3D NS): relative FD error vs.\ perturbation size~$\varepsilon$ at each rollout length~$T$.}
  \label{fig:app:ucurves_f3}
\end{figure}

\subsection{Optimization}
\label{app:recovery}

\paragraph{Optimizer configurations}
All optimization runs use either Adam (\texttt{optax.adam}, default
$\beta_1{=}0.9$, $\beta_2{=}0.999$, $\epsilon{=}10^{-8}$) or L-BFGS
(\texttt{optax.lbfgs} with the default zoom line search and a 10-step
inverse-Hessian memory), with patience-based early stopping on the loss:
the run terminates as soon as the loss has not strictly improved over
the running minimum for \emph{patience} consecutive iterations, or when
the maximum iteration budget is reached.
The 3D NS recovery and thermal conductivity convergence plots use
\emph{gradient evaluations} on the x-axis rather than outer iterations:
Adam costs one value-and-gradient call per step, while L-BFGS's zoom line
search adds an empirically observed average of $\sim$3 probes per outer
step, so a per-iteration x-axis would understate L-BFGS's true
forward-and-backward cost.
Per-task settings:
\begin{itemize}\itemsep0pt
  \item \textbf{H (conductivity recovery)}: Adam $\eta{=}10^{-2}$,
        2\,000 iters, patience 200; L-BFGS, 200 iters, patience 30.
  \item \textbf{S (topology optimization)}: Adam $\eta{=}5{\times}10^{-2}$,
        2\,500 iters, patience 100, with a soft volume-fraction penalty
        of weight 50; L-BFGS, 100 iters, patience 20 (fails to reach Adam's optimum, see below);
        MMA via NLopt~\citep{svanberg1987mma}, 200 iters, patience 30,
        with a hard volume inequality and native box constraints.
  \item \textbf{F2 (drag minimization, Re=20)}: Adam $\eta{=}5{\times}10^{-4}$,
        500 iters, patience 100, with a flow-rate penalty of weight 50;
        L-BFGS, 50 iters, patience 15 (diverges, see below).
  \item \textbf{F3 (3D NS IC recovery)}: Adam $\eta{=}10^{-3}$, 500 iters,
        patience 50; L-BFGS, 100 iters, patience 20.
        L-BFGS+proj and Adam+proj variants apply a solenoidal projection to
        the gradient at each step (see next paragraph) with otherwise
        identical settings. All variants start from $u\equiv 0$ and minimize
        the $\ell_2$ distance to the target final state, averaged over three
        IC seeds. The seeds vary the \emph{ground-truth target field}
        (\texttt{rand\_div\_free} regenerated with seeds $0,1,2$), not the
        optimizer initialization, which is identically zero across seeds;
        this measures recovery of arbitrary div-free targets from a cold
        start rather than re-running the same problem three times.
\end{itemize}

\paragraph{Solenoidal gradient projection (F3 variants)}
The L-BFGS+proj and Adam+proj variants on F3 apply a spectral
Helmholtz projection to the gradient \emph{before} each optimizer
update,
$\hat g^{\mathrm{df}}_i = \hat g_i - k_i (k\!\cdot\!\hat g) / |k|^2$,
where $\hat g$ is the FFT of the gradient on the periodic box. This
restricts the search direction to the divergence-free subspace and is
exact to machine precision on the periodic NS state space. The
projection is implemented once in the harness and applied uniformly to
all six fluid solvers; we do not project elsewhere because the
divergence-free constraint is specific to the incompressible NS state.

\paragraph{2D drag minimization}
\Cref{fig:app:dragopt} shows convergence for the 2D cylinder-flow drag minimization task.
This task is notably difficult to deploy: the channel geometry with advective outflow requires non-periodic pressure boundary conditions that exclude spectral solvers, while Brinkman penalization of the cylinder obstacle is incompatible with LU-factored Poisson projections, collectively ruling out four of the seven F2 backends (detailed in \cref{app:solver-issues}).
XLB and PICT converge to consistent drag reduction of around 60\% at Re\,=\,20; PhiFlow partially converges, reducing drag by ${\sim}35\%$ before stalling.
We also attempted L-BFGS on this task.
It failed universally: XLB diverged to NaN within 50 iterations (drag coefficients blowing up from 0.083 to over 8) and PICT produced no converged result.
The failure is structural: the inflow-to-drag map is strongly non-convex (advection-driven, with separation and re-attachment regimes), so secant pairs $s_k = x_{k+1} - x_k$ taken across iterations sample regions where the local Hessian differs in sign, violating the positive-curvature condition and causing the limited-memory model to blow up.
This contrasts with the quasi-static problems in \cref{fig:app:topopt,fig:app:conductivity_recovery}, where L-BFGS converges reliably; the determining factor is whether the landscape is stationary enough to sustain a coherent curvature model across iterations.

\begin{figure}[!htbp]
  \centering
  \includegraphics[width=\linewidth]{figures/drag_opt_re20.pdf}
  \caption{2D drag minimization at Re\,=\,20. XLB and PICT converge to ${\sim}60\%$ drag reduction; PhiFlow stalls at ${\sim}35\%$.}
  \label{fig:app:dragopt}
\end{figure}

\paragraph{3D NS IC recovery}
\Cref{fig:app:recovery} compares three optimizer configurations on the 3D NS IC recovery task ($[0,2\pi)^3$ TGV box, $16^3$ grid).
Each run starts from a constant zero-velocity field and minimizes the $\ell_2$ distance to a target final state.
We report the \emph{IC error} as the relative $\ell_2$ distance between the recovered and true initial velocity fields,
$\|\mathbf{u}^{*}_0 - \mathbf{u}_0\|_2 / \|\mathbf{u}_0\|_2$,
where $\mathbf{u}_0$ is the ground-truth IC and $\mathbf{u}^{*}_0$ is the optimizer iterate at the reported checkpoint.
Solver color and marker encode the backend; line style encodes the optimizer.
Adam stagnates above 40\% normalized IC error across all solvers, while L-BFGS converges below 6\% in fewer gradient evaluations.
Adding a solenoidal projection after each L-BFGS update (L-BFGS+proj) reduces the final IC error to below 0.5\% and keeps the recovered IC divergence at the physical level of the true IC throughout.
The projection benefit is specific to the second-order optimizer, and the asymmetry has a structural cause.
L-BFGS builds its Hessian approximation from secant pairs $s_k = x_{k+1} - x_k$; without projection these steps carry large irrotational components orthogonal to the solenoidal manifold, corrupting the curvature model, whereas projection confines $s_k$ to the correct tangent space and recovers the superlinear convergence behavior associated with projected Newton-type methods~\citep{bertsekas1982projected}.
Adam's per-coordinate second-moment scaling accumulates variance estimates from both solenoidal and irrotational gradient directions and has no cross-coordinate curvature model to exploit the constraint geometry~\citep{kingma2015adam}, so projecting the iterate does not repair the miscalibrated scaling; \cref{fig:app:recovery:adamproj} confirms this.
The bottom row shows $u_x$ at the middle $z$-slice for a representative solver (L-BFGS+proj/PhiFlow).

\begin{figure}[!htbp]
  \centering
  \includegraphics[width=\linewidth]{figures/recovery_overview.pdf}
  \caption{3D NS IC recovery overview for Adam, L-BFGS, and L-BFGS+proj (Adam+proj omitted; see \cref{fig:app:recovery:adamproj}).
    \emph{Top:} normalized IC error, IC divergence, and optimization loss vs.\ gradient evaluations across all solvers (color = solver, line style = optimizer).
    The dashed line marks the true IC divergence level.
    \emph{Bottom:} $u_x$ slice at convergence (L-BFGS+proj / PhiFlow): recovered IC, true IC, recovered final state, and true final state.}
  \label{fig:app:recovery}
\end{figure}

\begin{figure}[!htbp]
  \centering
  \includegraphics[width=\linewidth]{figures/recovery_adam_proj.pdf}
  \caption{Adam vs.\ Adam+proj on the 3D NS IC recovery task (linear scale).
    Adding a solenoidal projection step after each Adam update leaves IC error, IC divergence, and optimization loss essentially unchanged across all solvers, confirming that the projection benefit observed in \cref{fig:app:recovery} is specific to L-BFGS.}
  \label{fig:app:recovery:adamproj}
\end{figure}

\paragraph{Topology optimization}
\Cref{fig:app:topopt} summarizes the structural topology optimization task ($[0,2]\times[0,1]\times[0,1]$ cantilever beam, $16\times2\times8$ mesh, SIMP with $p=3$, Adam).
All solvers converge to the same final compliance (${\approx}0.001$): for smooth, well-conditioned physics, gradient quality does not limit solution quality.
The optimized density fields (bottom row) are visually identical across solvers, recovering the classic cantilever-beam truss topology.

We also evaluated L-BFGS and MMA~\citep{svanberg1987mma} on this task.
After 100 quasi-Newton steps, L-BFGS final compliance ranged from $0.001$ (TopOpt.jl) to $0.008$ (JAX-FEM), well above the Adam solution.
L-BFGS fails for two reinforcing reasons: the SIMP penalization ($p=3$) makes the objective globally non-convex, so the positive-curvature condition $y_k^\top s_k > 0$ required for valid BFGS updates is frequently violated and updates are skipped~\citep{fu2023quasinewton}; and box constraints $\rho \in [0,1]$ cause rapid active-set changes that invalidate the limited curvature history.
MMA, by contrast, converges to compliance ${\approx}0.00099$ in only 40 iterations, matching Adam's solution in under 2\% of its iteration budget, because it builds a strictly convex separable approximation at each step that sidesteps the SIMP non-convexity and handles box constraints and the volume inequality natively.
MMA is therefore the recommended optimizer for this task class, consistent with standard practice in the topology optimization community.

\begin{figure}[!htbp]
  \centering
  \includegraphics[width=\linewidth]{figures/topopt_overview.pdf}
  \caption{Structural topology optimization overview.
    \emph{Left:} final optimized density fields from a finer-mesh ($32{\times}4{\times}16$) MMA run, both as 3D voxel renderings (one per solver) and projected along each principal axis.
    \emph{Right:} compliance and volume-fraction vs.\ iteration for Adam and MMA (line style = optimizer, color = solver) on the coarser $16{\times}2{\times}8$ mesh used throughout the rest of the comparison; MMA converges in $\sim$40 iterations, Adam requires $\sim$2500.
    The two columns therefore use different mesh resolutions, as flagged in the panel headers: the finer run sharpens the recovered topology visually, while the coarser run keeps optimizer convergence directly comparable to the other tasks.}
  \label{fig:app:topopt}
\end{figure}

\paragraph{Conductivity recovery}
\Cref{fig:app:conductivity_recovery} shows convergence for the thermal conductivity inversion task.
The ground-truth conductivity field has two Gaussian peaks; the solver observes the steady-state temperature distribution under a known heat source and recovers the conductivity by minimizing the $\ell_2$ error between simulated and observed temperature.
Adam and L-BFGS are compared across solvers that support differentiation w.r.t.\ the conductivity field.
FEniCS and Firedrake form a distinct cluster: with Adam their final error is modestly higher ($\sim\!1{,}600$ vs.\ $\sim\!1{,}300$), and with L-BFGS the gap widens sharply ($\sim\!8.5{\times}10^4$ vs.\ $\sim\!1{,}100$).
The cause is a gradient inconsistency in the adjoint-based solvers: the forward pass reports $\sum_i(T_i - T_i^*)^2$ (a nodal sum), while the adjoint differentiates $\int(T-T^*)^2\,\mathrm{d}\Omega$ (an area-weighted integral) and applies a scalar correction factor $n_\text{nodes}/|\Omega|$ to bridge the two.
This correction is only approximate, so the returned gradient is not the exact derivative of the reported objective.
First-order methods tolerate this mismatch because they use only gradient direction; L-BFGS accumulates gradient differences to estimate curvature, so an inconsistent gradient corrupts the Hessian approximation and causes convergence to a suboptimal point.

\begin{figure}[!htbp]
  \centering
  \includegraphics[width=\linewidth]{figures/conductivity_recovery_overview.pdf}
  \caption{Thermal conductivity recovery comparing Adam and L-BFGS optimizers. \emph{Top:} identification error vs.\ gradient evaluations; color encodes solver, linestyle encodes optimizer. \emph{Bottom:} final recovered conductivity profiles (all solvers) vs.\ ground truth (dashed).}
  \label{fig:app:conductivity_recovery}
\end{figure}


\section{Solver documentation}
\label{app:solvers}

This appendix documents all solvers evaluated for inclusion in Mosaic,
covering implementation details, wrapping approach, and known limitations.
Solvers included in the benchmark appear in Table~\ref{tab:solvers}.
Excluded solvers and the reasons for their exclusion are listed in
\cref{app:excluded}.

\subsection{Included solvers}
\label{app:included}

\paragraph{JAX-CFD~\cite{kochkov2021machine}}
JAX-native incompressible flow solver on a staggered MAC grid with
finite-difference advection and spectral (FFT) pressure projection.
Gradients via JAX source-transformation AD.
The spectral pressure solve requires periodic BCs in all spatial directions,
so JAX-CFD is included in periodic benchmarks (TGV, multimode
agreement) but excluded from the cylinder-wake experiment.

\paragraph{PhiFlow~\cite{holl2024phiflow}}
Semi-Lagrangian advection with pressure projection.
Supports PyTorch, JAX, and TensorFlow backends, enabling gradient
computation through the same simulation code across AD frameworks.

\paragraph{INS.jl~\cite{agdestein2024ins}}
IncompressibleNavierStokes.jl: Julia finite-difference pressure-projection
solver. Differentiates through the time loop via Zygote.jl reverse-mode AD.
CPU only.

\paragraph{XLB~\cite{ataei2024xlb}}
JAX-native lattice Boltzmann solver supporting D2Q9 (2D) and D3Q27 (3D) stencils on
GPU. Gradients via source-transformation AD through the collision-streaming
loop.

\paragraph{PICT~\cite{franz2025pict}}
GPU-accelerated differentiable incompressible Navier-Stokes solver built on
PyTorch with custom CUDA kernels implementing the PISO algorithm. Supports
reverse-mode AD through the full time loop and handles multi-block
curvilinear grids. The Mosaic Tesseract exposes periodic, lid-driven cavity,
inflow/channel (with advective outflow), and cylinder-wake modes via an
8-block ring topology.

\paragraph{Warp-NS~\cite{macklin2024warp}}
Custom incompressible NS solver built with NVIDIA Warp CUDA kernels,
implementing the IPCS projection scheme with an FFT-based Poisson solve
(following a reference spectral example; a non-periodic solver was not
implemented). Gradients via the \texttt{wp.Tape} kernel-level VJP.
The Warp framework provides differentiable kernels but no ready-made
incompressible flow solver, so the FD stencils, pressure Poisson solve,
and time integration were implemented from scratch using Warp primitives.
Warp-NS represents a best-effort implementation by the Mosaic team and
illustrates the integration cost when a kernel toolkit provides no
built-in solver.

\paragraph{Exponax~\cite{koehler2024apebench}}
Spectral incompressible Navier-Stokes solver supporting both 2D
(streamfunction-vorticity) and 3D (Leray-projected velocity) formulations,
integrated with an exponential time-differencing Runge-Kutta (ETDRK) scheme.
Incompressibility is enforced to machine precision by construction.
Gradients via JAX source-transformation AD.
Used in Mosaic for the 3D periodic domain (F3).

\paragraph{FEniCS~\cite{farrell2013automated}}
Finite-element solver using P1 elements for thermal and structural problems.
dolfin-adjoint automates the discrete adjoint by replaying the forward tape.

\paragraph{Firedrake~\cite{rathgeber2016firedrake}}
Mirrors the FEniCS P1/CG1 formulation for structural and thermal problems.
Differentiates via firedrake-adjoint, providing an independent
tape-based adjoint implementation for cross-validation.

\paragraph{JAX-FEM~\cite{xue2023jaxfem}}
Solves heat conduction and linear elasticity with trilinear HEX8 finite
elements in JAX. Gradients via AD through the assembled system.

\paragraph{TopOpt.jl~\cite{huang2021topoptjl}}
SIMP topology optimization for linear elasticity with HEX8 elements in
Julia, using analytical adjoint sensitivities.

\paragraph{torch-fem~\cite{meyer2024torchfem}}
PyTorch finite-element solver for heat conduction with linear HEX8 elements.
Gradients via PyTorch autograd through the assembled system.
GPU-accelerated via PyTorch sparse operations.

\paragraph{OpenFOAM~\cite{weller1998tensorial} (reference)}
Runs the icoFoam incompressible PISO solver as a forward-only reference
baseline. No reverse-mode AD available.

\paragraph{deal.II~\cite{bangerth2007dealii,arndt2023dealii} (reference)}
Solves thermal and structural problems with Q1 elements using the
industry-grade C++ finite-element library.
No native reverse-mode AD is available; used as a forward-only reference
baseline for S and H, analogously to OpenFOAM for the fluid domains.

\subsection{Excluded solvers}
\label{app:excluded}

We document solvers that were evaluated for inclusion but ultimately
excluded, along with the specific reason in each case. Practitioners
looking for a particular solver can check whether it was considered and
why it was left out.

\paragraph{WaterLily.jl~\cite{waterlily2024}}
Forward-mode AD only, via \texttt{ForwardDiff.jl} dual numbers. No
Zygote reverse-mode support is documented. A VJP can be emulated by
contracting $N$ forward-mode JVPs, but the cost scales as
$O(N^d \cdot T_{\mathrm{fwd}})$, matching central finite differences.
Excluded from gradient benchmarks. Forward accuracy is reported where
applicable.

\paragraph{JAX-Fluids~\cite{jaxfluids2022}}
Targets compressible flow. Recovering the incompressible limit at low Mach
number ($\mathrm{Ma} \approx 0.01$) requires acoustic sub-stepping whose
iteration count depends on Ma, grid size, and time step simultaneously.
Because \texttt{jax.lax.scan} requires a statically known loop count, a
fully differentiable low-Mach wrapper cannot be built around the published
library without reimplementing its time integrator. Excluded from both
forward and gradient benchmarks.

\paragraph{Commercial solvers (ANSYS, COMSOL, Abaqus)}
Excluded for reproducibility: closed-source licenses prevent redistribution
of solver binaries or independent verification of gradient implementations.

\paragraph{FEniCS-NS~\cite{farrell2013automated}}
The FEniCS incompressible NS solver (CG1/CG1 elements, IPCS time-stepping)
was evaluated for fluid benchmarks but excluded for two independent reasons.
Without SUPG/GLS advection stabilization, stability requires
$\mathrm{Pe}_h = Uh/(2\nu) \lesssim 1$; at $\nu < 0.01$ on the $N=64$
benchmark grid, relative errors reach 31--212 while other solvers stay in
0.02--0.16.
Independently, wall-clock scaling is empirically $O(N^{2.45})$ in serial;
the VJP at $N=64$ takes 1248\,s, and the forward pass alone for the drag
optimization task ($N=32$, 400 steps) requires approximately 25{,}000\,s,
exceeding the 1200\,s HTTP watchdog by a factor of 40.
FEniCS remains included for structural and heat-transfer domains (H, S),
where neither constraint applies.

\subsection{Observed solver limitations}
\label{app:solver-issues}

The following limitations were identified during benchmarking and resulted
in task-specific exclusions. Solvers excluded globally are documented in
\cref{app:excluded}.

\paragraph{PICT: viscosity not differentiable in PISOtorch\_diff}
PISOtorch\_diff tracks autograd through velocity fields and boundary
conditions, but treats viscosity as a static scalar: passing a
\texttt{requires\_grad=True} tensor as the viscosity argument produces a
result with no \texttt{grad\_fn} (confirmed by direct inspection of the
autograd graph). Differentiation w.r.t.\ $\nu$ would require adding
viscosity as an explicit differentiable parameter inside the PISOtorch\_diff
C++/CUDA kernels; this is outside the scope of the Mosaic wrapper.

\paragraph{JAX-CFD: spectral pressure solve requires periodic boundary conditions}
JAX-CFD uses an FFT-based pressure Poisson solver that requires doubly-periodic
boundary conditions by construction. The 2D cylinder drag optimization
domain is a channel with inlet, outlet, and obstacle boundaries, none of
which are periodic. Volume penalization masks velocity inside the solid
region but does not remove the periodicity requirement from the pressure
solve. JAX-CFD is structurally excluded from the drag optimization task.

\paragraph{Warp-NS: no non-periodic pressure solver available within the time budget}
The Warp-NS implementation was built following a reference example that uses an
FFT-based spectral Poisson solver, which inherits the same doubly-periodic BC
requirement. Unlike JAX-CFD, this is not a structural limitation of the Warp
framework: a non-periodic iterative Poisson solver could in principle be
implemented using Warp primitives. However, no suitable example was available,
and implementing one from scratch was outside the time budget of this work;
the drag optimization was therefore not attempted.

\paragraph{INS.jl: viscosity gradient not available}
Zygote.jl cannot differentiate through the diffusion term in
IncompressibleNavierStokes.jl because its \texttt{ChainRulesCore} rule for
the diffusion operator returns \texttt{NoTangent()} for the viscosity
argument. Gradients w.r.t.\ $\nu$ are therefore not reported for INS.jl.
Fixing this requires either upstreaming a correct rrule into INS.jl or
implementing a custom adjoint that accounts for the diffusion term.

\paragraph{INS.jl: Brinkman penalization incompatible with spectral pressure solve}
INS.jl (IncompressibleNavierStokes.jl)~\cite{agdestein2024ins} supports periodic,
Dirichlet, symmetric, and advective outlet boundary conditions, but not
immersed boundary (IBM) or cut-cell methods. Applying Brinkman volume
penalization to represent the cylinder obstacle introduces a velocity
discontinuity at the solid boundary after each spectral LU pressure
projection step; the Poisson solve amplifies this discontinuity, and the
velocity field diverges to NaN at all tested resolutions ($N \geq 16$). A
correct approach requires incorporating the penalization term directly into
the pressure Poisson system~\cite{angot1999penalization}, which is outside
the scope of the current solver implementation. INS.jl is excluded from the
drag optimization task.

\paragraph{XLB: intrinsic $O(\mathrm{Ma}^2)$ compressibility error}
The D2Q9 BGK lattice Boltzmann method~\cite{qian1992lattice} recovers the
incompressible Navier-Stokes equations only to $O(\mathrm{Ma}^2)$ via
Chapman-Enskog expansion. At fixed $\Delta t$, the Mach number
$\mathrm{Ma} = u\,\Delta t / \Delta x$ grows with grid refinement, so the
compressibility error does not decrease with spatial resolution for
incompressible problems. At $N = 128$, $\Delta t = 0.01$,
$\mathrm{Ma} \approx 0.2$, giving an $O(0.04)$ error floor independent of
solver correctness~\cite{ataei2024xlb}. XLB forward accuracy results at
fine grids should be interpreted with this limit in mind.

\paragraph{XLB: BGK collision instability at low viscosity}
The BGK relaxation time $\tau = \nu/c_s^2\,\Delta t + 0.5$ approaches the
stability boundary $\tau \to 0.5$ as $\nu \to 0$~\cite{qian1992lattice}.
The XLB tesseract automatically selects the KBC (entropic) collision
operator~\cite{ataei2024xlb} when $\omega > 1.8$ ($\tau < 0.556$), which is
unconditionally entropy-stable.
Results at $\nu \in \{10^{-4},\, 5\times10^{-4}\}$ are pending re-validation
under the KBC operator and should be interpreted with caution.


\end{document}
